\thispagestyle{normalpage}

\begin{center}
    \large \textbf{CONCLUSION GÉNÉRALE}
\end{center}
\addcontentsline{toc}{chapter}{CONCLUSION GÉNÉRALE}
\label{page:conclusion}

\vspace{0.5cm}

Ce mémoire avait pour objectif de répondre à une problématique centrale : comment faire évoluer une plateforme pharmaceutique comme Techcare, confrontée aux limites structurelles d'une architecture monolithique, afin de garantir la scalabilité, la résilience et la maintenabilité requises par sa croissance ?

\vspace{0.3cm}

Pour y répondre, nous avons adopté une démarche scientifique et progressive. Le premier chapitre a posé le contexte du projet et les exigences fonctionnelles et non-fonctionnelles. Le deuxième chapitre a fourni les fondements théoriques indispensables à toute architecture distribuée — le théorème CAP, les patterns de résilience, le Domain-Driven Design — et a abouti à une conception rigoureuse de l'architecture cible. Le troisième chapitre a concrétisé cette conception en une implémentation technique avec NestJS, RabbitMQ, Docker et les services cloud AWS. Enfin, le quatrième chapitre a offert un regard critique sur les résultats obtenus et les perspectives d'évolution.

\vspace{0.3cm}

Les objectifs fixés au début de ce projet ont été pleinement atteints. Sur le plan de la scalabilité, l'architecture microservices permet désormais de faire évoluer indépendamment chaque domaine métier. Par exemple, le service de gestion des produits peut absorber des pics de consultation importants sans que cela n'impacte le traitement des paiements. En matière de résilience, l'intégration du pattern Saga avec un mécanisme de compensation garantit l'intégrité des données distribuées, assurant ainsi que la panne d'un service n'entraîne plus l'indisponibilité totale de la plateforme. Concernant la maintenabilité, l'adoption d'une architecture hexagonale au sein de chaque service, couplée à l'utilisation du langage ubiquiste issu du Domain-Driven Design, a permis de concevoir un code parfaitement aligné sur le métier pharmaceutique. Cela réduit considérablement les risques de régression lors des évolutions futures. Enfin, du point de vue de la sécurité, la mise en place d'une stratégie Zero Trust, la propagation de l'identité via des jetons JWT et un contrôle d'accès basé sur les rôles (RBAC) offrent un niveau de protection optimal et adapté aux exigences d'un contexte médical sensible.

\vspace{0.3cm}

Néanmoins, cette transition n'est pas sans limites. La complexité opérationnelle accrue, la gestion de la consistance éventuelle, ainsi que le léger compromis sur la latence réseau (le fameux "impôt de la distribution") représentent des défis concrets inhérents aux systèmes distribués. Ces limites, identifiées et mesurées lors de nos phases de test, constituent une véritable feuille de route pour les évolutions futures. Les perspectives d'optimisation sont claires : l'orchestration dynamique avec Kubernetes, l'adoption d'un \textit{Service Mesh} pour déléguer la sécurité inter-services et renforcer l'observabilité existante sans modifier le code applicatif, ou encore l'intégration d'un cache Redis et du protocole gRPC pour pulvériser les temps de réponse.

\vspace{0.3cm}

En définitive, ce projet démontre qu'une migration vers les microservices ne doit pas être menée comme une fin en soi, mais comme une réponse mesurée à des besoins métiers concrets et quantifiables. 
